\documentclass[10pt,a4paper]{amsart}
\usepackage[margin=19mm]{geometry}
\usepackage{amsmath,amssymb,mathtools}
\usepackage[scr=rsfs,cal=euler]{mathalfa}
\usepackage{xfrac}
\usepackage[unicode,hidelinks,bookmarks=false]{hyperref}
\newcommand{\ie}{i.e.\ }
\newcommand{\cf}{cf.\ }
\newcommand{\resp}{respectively\ }
\newcommand{\Subsection}{\subsection}
\providecommand{\nobreakdash}{\nobreak\hbox{-}}
\setlength{\emergencystretch}{2em}
\multlinegap0pt
\usepackage{bm}
\usepackage{amssymb}
\usepackage[shortlabels]{enumitem}
\usepackage{mathtools}
\usepackage{algorithm}\usepackage{algpseudocode}
\newtheorem{theorem}{Theorem}
\def\E{\mbb{E}} %
\def\KL#1#2{\textnormal{KL}({#1}\Vert{#2})}
\newcommand{\RR}{\mathbb{R}}
\newcommand{\df}{\mathrm{d}}
\def\mbb#1{\mathbb{#1}}
\title{Probabilistic Inference and Learning with Stein's Method}
\author{Qiang Liu, Lester Mackey, Chris Oates}
\date{Source: arXiv:2603.07467v1 (CC BY 4.0)}
\begin{document}
\maketitle

\noindent\emph{Source and licence.} Probabilistic Inference and Learning with Stein's Method. Version arXiv:2603.07467v1 (CC BY 4.0), distributed by arXiv under the Creative Commons Attribution 4.0 International licence, http://creativecommons.org/licenses/by/4.0/, as recorded in the arXiv licence metadata for this exact version and shown on the abstract page https://arxiv.org/abs/2603.07467v1. Full licence notice: \texttt{licenses/arxiv-2603.07467-CC-BY-4.0.txt}.\par\smallskip

\noindent\textit{Editorial scope.} Counted unit: the article's theorem KL (source0.tex lines 4077--4092). The article's complete original proof of it is delivered with it (source0.tex lines 4094--4131, one \texttt{proof} environment of 38 lines). No mathematics is written, repaired or re-derived: the statement and the proof are the article's own lines, in the article's own notation, and no retained line is substituted or re-flowed.  Retained uncounted as the source-local context the proof needs: lines 4060--4063.  Nothing was omitted from any retained span, so the package carries no omission record.  No retained line contains a citation, so the wrapper carries no bibliography and no bibliography entry.  No retained line contains a figure environment, an image-inclusion command or drawing code, so no image is delivered and the package declares no asset dependency.  Editorial formatting only, in addition: an AMS \texttt{amsart} preamble that restates the article's own declarations and macros used by the retained lines, namely the environments \texttt{theorem} on separate counters, together with the standard packages those lines need.  The retained environments print in this document as Theorem 1 in order of appearance, whereas the article prints the same items with the numbering of its own full document.  The complete native source of the article as distributed in the arXiv e-print for this exact version is bundled unchanged as \texttt{sources/195944/source0.tex}.
\begin{align}
\label{equ:liouville}
\frac{\df}{\df t} \log q_t(Z_t) = - (\nabla \!\cdot v_t)(Z_t).
\end{align}

\begin{theorem}\label{thm: deriv KL}
Assume that the differential equation $\dot Z_t = v_t(Z_t)$ admits a unique solution on $[0,T]$ for every initialization in $\RR^d$.
Let $Q_t$ denote the law of $Z_t$, with density $q_t$.
Assume that $v_t$ and $\log q_t$ are continuously differentiable on $[0,T] \times \RR^d$.
Let $P$ be a probability measure with a positive and continuously differentiable density $p$, such that $\KL{Q_0}{P} < \infty$, and 
$\int_0^T Q_t(|(\mathcal{T}_P v_t)(Z_t)|)\df t < \infty,$ 
where $\mathcal{T}_P$ is the Langevin Stein operator for $P$.
Then
\begin{align}
\frac{\df}{\df t}\,
\KL{Q_t}{P}
=
-%
\,Q_t(\mathcal T_P v_t). 
\end{align} 
\end{theorem}

\begin{proof}  
Since $Z_t$ follows law $Q_t$, we have  
\[
\KL{Q_t}{P} = \E\left[\log q_t(Z_t) - \log p(Z_t)\right]. 
\]
Define $R_t = \log q_t(Z_t) - \log p(Z_t)$. Its time derivative is
\begin{align*}  
\dot R_t 
&= \frac{\df}{\df t} \log q_t(Z_t) - \frac{\df}{\df t} \log p(Z_t)  \\
&= - \nabla \cdot v_t(Z_t) - v_t(Z_t) \cdot \nabla \log  p(Z_t) \\
&= -\mathcal{T}_P v_t(Z_t),
\end{align*}
where we used Liouville’s formula \eqref{equ:liouville} for the first term,
and the chain rule for the second term,
\[
\frac{\df}{\df t} \log p(Z_t)
=
\nabla \log p(Z_t)^\top \dot Z_t
=
\nabla \log p(Z_t)^\top v_t(Z_t).
\]
Therefore,
\begin{align*}
\KL{Q_t}{P} - \KL{Q_0}{P} 
&= \E\left[R_t - R_0\right] \\
&= \E\left[\int_0^t \dot R_\tau\, \df \tau \right] \\
&= \int_0^t \E\left[\dot R_\tau\right]\, \df \tau \\
&= -\int_0^t \E\left[(\mathcal{T}_P v_\tau)(Z_\tau)\right]\, \df \tau,
\end{align*}
where we used Fubini--Tonelli’s theorem to exchange the order of integration, assuming  
\(\int_0^t \E [ \lvert (\mathcal{T}_P v_\tau)(Z_\tau)\rvert ]\, \df \tau < \infty\).
Differentiating the resulting identity and using continuity of  
\( t \mapsto \E \left[\mathcal T_P v_t(Z_t)\right] \) yields
$ 
\frac{\mathrm{d}}{\mathrm{d}t}\KL{Q_t}{P}
 = -\,\E\left[(\mathcal{T}_P v_t)(Z_t)\right],$
concluding the proof.
\end{proof}

\end{document}
